\section*{Capítulo \ref{ch:b2:ore-deposits} --- Yacimientos minerales}

\begin{solution}{exo:b2:ore-deposits:1}
(a) $\qty{0.5}{\percent} = \qty{5000}{ppm}$, y $5000/28 = 179$. (b) $\qty{6}{g/t} = \qty{6000}{ng/g}$,
y $6000/1.5 = 4000$. (c) La cromita contiene una fracción másica de cromo de $2 \times 52.00/(55.85 + 104.00 + 64.00) = 0.4646$;
la cromitita, $0.9 \times 0.4646 = 0.418$, es decir, \qty{418000}{ppm}, y
$418000/92 = 4540$.
\end{solution}

\begin{solution}{exo:b2:ore-deposits:2}
(a) $914 \times 0.0053 = \qty{4.84}{Mt}$ de cobre y $914 \times 10^6 \times 0.40 = \num{3.66e8}$ g,
\qty{366}{t} de oro; $1640 \times 0.0044 = \qty{7.22}{Mt}$ y \qty{738}{t}. (b) $1640 - 914 =
\qty{726}{Mt}$ que contienen $7.216 - 4.844 = \qty{2.37}{Mt}$ de cobre, con una \omterm{def:b2:ore-deposits:grade}{ley} media de
$2.372/726 = \qty{0.327}{\percent}$: los bloques adicionales están por debajo de la \omterm{def:b2:ore-deposits:grade}{ley de corte} del \qty{0.3}{\percent}
y justo por encima de ella. (c) Responden a la pregunta para dos \omterm{def:b2:ore-deposits:grade}{leyes de corte}; ninguna es el tamaño del
yacimiento sin su \omterm{def:b2:ore-deposits:grade}{ley de corte}.
\end{solution}

\begin{solution}{exo:b2:ore-deposits:3}
(a) $(19.3 - 1.00)/(2.65 - 1.00) = 18.3/1.65 = 11.1$. (b) Velocidades iguales exigen $(\rho_g - 1)d_g^2 =
(\rho_q - 1)d_q^2$, así que $d_g/d_q = \sqrt{1.65/18.3} = 0.300$: un grano de oro sedimenta como un grano de cuarzo
unas tres veces mayor. (c) $d_q = 0.1\sqrt{3.7/1.65} = \qty{0.150}{mm}$.
\end{solution}

\begin{solution}{exo:b2:ore-deposits:4}
Por unidad de superficie y de densidad, la capa conserva su propio cobre y recibe el de la columna
lixiviada: $g = g_0(h_e + \eta h_l)/h_e$. (a) $0.4 \times (50 + 150)/50 = \qty{1.6}{\percent}$.
(b) $0.4 \times (50 + 0.8 \times 150)/50 = \qty{1.36}{\percent}$. (c) $h_l = 50 \times (1.7/0.4 - 1)
= \qty{163}{m}$. Estas \omterm{def:b2:ore-deposits:grade}{leyes} caen en el intervalo 0.4--\qty{1.7}{\percent} de las capas de enriquecimiento
de los pórfidos.
\end{solution}

\begin{solution}{exo:b2:ore-deposits:5}
(a) Platino, $C_S = 10 \times 10^4 (R + 1)/(R + 10^4)$ ppb: \qty{1.00}{ppm} para $R = 100$,
\qty{9.10}{ppm} para \num{1000}, \qty{50.0}{ppm} para \num{10000}, \qty{90.9}{ppm} para \num{100000}.
Níquel, $C_S = 200 \times 500 (R + 1)/(R + 500)$ ppm: \qty{1.68}{\percent}, \qty{6.67}{\percent},
\qty{9.52}{\percent}, \qty{9.95}{\percent}. (b) $D\,C_L$: \qty{100}{ppm} de platino,
\qty{10}{\percent} de níquel. (c) El níquel está a menos del \qty{5}{\percent} de su techo a partir de
$R = \num{10000}$, así que su contenido apenas cambia en un factor diez de $R$; el platino todavía
casi se duplica en ese intervalo, de modo que su contenido mide $R$.
\end{solution}

\begin{solution}{exo:b2:ore-deposits:6}
(a) $\mu_M = 2.373 - 0.364 - 2 = 0.009$; $\sigma_M = 0.702$; mediana $10^{0.009} = \qty{1.02}{Mt}$.
(b) $z = (1 - 0.009)/0.702 = 1.412$, $1 - \Phi(1.412) = 0.0790$: unos 33 yacimientos de 422.
(c) $\sigma = \sqrt{0.654^2 + 0.199^2} = 0.684$, $z = 1.449$, $1 - \Phi = 0.0737$, unos 31: la
débil correlación añade unos pocos gigantes.
\end{solution}

\begin{solution}{exo:b2:ore-deposits:7}
Media de todos los bloques $0.40\,e^{0.18} = \qty{0.479}{\percent}$. Para $c = 0.3$: $z = \ln(0.40/0.3)/0.60 =
0.4795$, $\Phi(z) = 0.6842$, $\Phi(z + s) = 0.8598$: $T = \qty{684}{Mt}$, $\bar g = 0.479 \times
0.8598/0.6842 = \qty{0.602}{\percent}$, $M = 1000 \times 0.00479 \times 0.8598 = \qty{4.12}{Mt}$.
Para $c = 0.6$: $z = -0.6758$, $\Phi = 0.2496$, $\Phi(z + s) = 0.4698$: \qty{250}{Mt} al
\qty{0.901}{\percent}, \qty{2.25}{Mt}. En la figura, para una \omterm{def:b2:ore-deposits:grade}{ley de corte} del \qty{0.3}{\percent} la curva azul
marca unas \qty{680}{Mt} y la roja alrededor del \qty{0.6}{\percent}.
\end{solution}

\begin{solution}{exo:b2:ore-deposits:8}
(a) Fluido: $\num{5e6}/(\num{1e-3} \times 0.5) = \qty{1.00e10}{t}$; \omterm{def:b2:magma-emplacement:magma}{magma}: $\num{1e10}/0.035 =
\qty{2.86e11}{t}$. (b) $\num{2.86e14}/2600 = \qty{1.10e11}{m^3}$, \qty{110}{km^3}, un cubo de \omterm{def:b2:magma-emplacement:magma}{magma}
de unos \qty{4.8}{km} de arista. (c) El \omterm{def:b2:magma-emplacement:magma}{magma} contenía $\num{2.86e11}
\times \num{50e-6} = \qty{1.43e7}{t}$ de cobre, y la mena $5/14.3 = 0.350$ de él.
\end{solution}

\begin{solution}{exo:b2:ore-deposits:9}
(a) Media de los logaritmos 1.411, desviación típica 0.0664; mediana $10^{1.411} = \qty{25.8}{ppm}$;
umbral $10^{1.411 + 0.133} = 10^{1.544} = \qty{35.0}{ppm}$. (b) $(4 \times 400 + 36 \times 25.8)/40 =
\qty{63.2}{ppm}$, por encima de \qty{35.0}{ppm}: sigue siendo anómalo tras la dilución con nueve partes de sedimento
limpio. (c) $4 \times 400 + (A - 4) \times 25.8 = 35.0\,A$ da $A = 4 \times (400 -
25.8)/(35.0 - 25.8) = \qty{163}{km^2}$.
\end{solution}

\begin{solution}{exo:b2:ore-deposits:10}
(a) $e^{su}\,e^{-u^2/2} = e^{-(u - s)^2/2 + s^2/2}$; divídase por $\sqrt{2\pi}$. (b) $E[e^{\mu + su}] =
e^{\mu}\int e^{su}\varphi(u)\,\mathrm{d}u = e^{\mu + s^2/2}\int\varphi(u - s)\,\mathrm{d}u = e^{\mu +
s^2/2}$. (c) $\Phi$ es creciente y $s > 0$, así que $\Phi(z + s) > \Phi(z)$: los bloques conservados son más ricos que
los desechados, de modo que cada \omterm{def:b2:ore-deposits:grade}{ley de corte} conserva una parte mayor del metal que de la roca; la
diferencia crece con $s$, la variabilidad del yacimiento.
\end{solution}

\begin{solution}{exo:b2:ore-deposits:11}
(a) $D(R + 1)/(R + D) = fD$ da $R + 1 = fR + fD$, $R = (fD - 1)/(1 - f)$. (b) $f = 0.5$: $R = D - 2$,
298 y \num{9998}; $f = 0.9$: $R = 9D - 10$, \num{2690} y \num{89990}. (c) El níquel alcanza el
\qty{50}{\percent} de su techo dentro del intervalo de las menas de níquel y cobre (de 100 a \num{2000}), y el 90 justo por encima;
el platino necesita $R$ de \num{10000} a \num{100000}, el intervalo de los niveles de EGP.
\end{solution}

\begin{solution}{exo:b2:ore-deposits:12}
(a) Con $\ln M = m + s u$, los yacimientos por encima del cuantil tienen $u > \Phi^{-1}(q)$; su metal es
$\int_{\Phi^{-1}(q)}^{\infty} e^{m + su}\varphi(u)\,\mathrm{d}u = e^{m + s^2/2}[1 - \Phi(\Phi^{-1}(q) -
s)]$, y dividir por el total $e^{m + s^2/2}$ da $\Phi(s - \Phi^{-1}(q))$. (b) $s = 0.702 \ln 10 =
1.616$; $\Phi(1.616 - 1.282) = \Phi(0.334) = 0.631$. (c) $\Phi(1.616 - 2.326) = \Phi(-0.710) = 0.239$,
frente a 0.266 para los cuatro mayores: los yacimientos más grandes son algo mayores de lo que espera la ley
lognormal.
\end{solution}

\begin{solution}{pb:b2:ore-deposits:1}
\textbf{1.} \qty{330}{m}; $\sum l\,g = 138.4$, media $138.4/330 = \qty{0.419}{\percent}$.
\textbf{2.} $3.24/8 = \qty{0.405}{\percent}$; cada tramo representa una longitud de testigo, y por tanto un
volumen de roca: un tramo largo debe contar más que uno corto.
\textbf{3.} $230/330 = 0.697$, con $118.6/230 = \qty{0.516}{\percent}$.
\textbf{4.} $4190/28 = 150$.
\textbf{5.} $8100/28 = 289$.
\textbf{6.} $0.38\,e^{0.55^2/2} = 0.38\,e^{0.1513} = \qty{0.442}{\percent}$; el sondeo, con un
\qty{0.419}{\percent}, queda a menos del \qty{6}{\percent} de ese valor: es representativo.
\textbf{7.} $\pi \times 450^2 \times 700 = \qty{4.45e8}{m^3}$; $\times 2.7 = \num{1.20e9}$ t,
\qty{1200}{Mt}.
\textbf{8.} $z = \ln(0.38/0.3)/0.55 = 0.4298$; $\Phi(z) = 0.6663$.
\textbf{9.} $z = 1.167$, 0.4298, $-0.4990$; $\Phi = 0.8784$, 0.6663, 0.3089; $T = \qty{1050}{Mt}$,
\qty{800}{Mt}, \qty{371}{Mt}.
\textbf{10.} $\Phi(z + s) = \Phi(0.9798) = 0.8364$; $\bar g = 0.442 \times 0.8364/0.6663 =
\qty{0.555}{\percent}$.
\textbf{11.} $0.442 \times 0.9570/0.8784 = \qty{0.482}{\percent}$; $0.442 \times 0.5203/0.3089 =
\qty{0.744}{\percent}$.
\textbf{12.} \omterm{def:b2:ore-deposits:grade}{Tonelaje} $\times\, 370.7/799.6 = 0.464$; \omterm{def:b2:ore-deposits:grade}{ley} $\times\, 0.744/0.555 = 1.34$.
\textbf{13.} En la mediana, \qty{0.38}{\percent} ($z = 0$): \qty{600}{Mt} con $0.442 \times
\Phi(0.55)/0.5 = 0.442 \times 0.7088/0.5 = \qty{0.627}{\percent}$.
\textbf{14.} $1200 \times 0.00442 = \qty{5.30}{Mt}$.
\textbf{15.} $(\log_{10} 1200 - 2.373)/0.654 = (3.0792 - 2.373)/0.654 = 1.080$; $\Phi = 0.8599$: el
percentil 86, un yacimiento grande.
\textbf{16.} $(\log_{10} 0.442 + 0.364)/0.199 = (-0.3546 + 0.364)/0.199 = 0.047$; $\Phi = 0.5187$: una
\omterm{def:b2:ore-deposits:grade}{ley} corriente, percentil 52.
\textbf{17.} 0.702 con la correlación, 0.684 sin ella.
\textbf{18.} $\mu_M = 0.009$; $(\log_{10} 5.30 - 0.009)/0.702 = (0.7243 - 0.009)/0.702 = 1.019$; $\Phi =
0.8459$: el percentil 85.
\textbf{19.} $422 \times (1 - 0.8459) = 65$.
\textbf{20.} $M = T_0\,e^{\mu + s^2/2}\,\Phi(z + s) = 1200 \times 0.00442 \times 0.9570 =
\qty{5.08}{Mt}$.
\textbf{21.} $1200 \times 0.00442 \times 0.5203 = \qty{2.76}{Mt}$.
\textbf{22.} $\Phi(z + s) = 0.8364$ del cobre, en $\Phi(z) = 0.6663$ del \omterm{def:b2:ore-deposits:grade}{tonelaje}.
\textbf{23.} Del 0.3 al \qty{0.2}{\percent}, el \omterm{def:b2:ore-deposits:grade}{tonelaje} crece en $1054 - 800 = \qty{254}{Mt}$
(\qty{32}{\percent}) y el cobre en $5.08 - 4.44 = \qty{0.64}{Mt}$ (\qty{14}{\percent}): los bloques añadidos
tienen una \omterm{def:b2:ore-deposits:grade}{ley} media de $0.64/254 = \qty{0.25}{\percent}$, los más pobres del yacimiento.
\textbf{24.} $1200 \times 0.00442 \times 0.8364 = 800 \times 0.00555$: el prospecto contiene
\textbf{\qty[text-series-to-math=true]{4.44}{Mt}} de cobre por encima de la \omterm{def:b2:ore-deposits:grade}{ley de corte} del \qty{0.3}{\percent}.
\end{solution}
